
\documentclass{article}
\usepackage{colortbl}
\usepackage{makecell}
\usepackage{multirow}
\usepackage{supertabular}

\begin{document}

\newcounter{utterance}

\centering \large Interaction Transcript for game `dond', experiment `semi\_en', episode 13 with reflexion{-}Qwen3.5{-}27B.
\vspace{24pt}

{ \footnotesize  \setcounter{utterance}{1}
\setlength{\tabcolsep}{0pt}
\begin{supertabular}{c@{$\;$}|p{.15\linewidth}@{}p{.15\linewidth}p{.15\linewidth}p{.15\linewidth}p{.15\linewidth}p{.15\linewidth}}
    \# & $\;$A & \multicolumn{4}{c}{Game Master} & $\;\:$B\\
    \hline

    \theutterance \stepcounter{utterance}  
    & & \multicolumn{4}{p{0.6\linewidth}}{
        \cellcolor[rgb]{0.9,0.9,0.9}{
            \makecell[{{p{\linewidth}}}]{
                \texttt{\tiny{[P1$\langle$GM]}}
                \texttt{You have attempted this same task before and did not succeed. These are the plans you wrote after those attempts.} \\
\texttt{Trial \#1: New plan: I failed because my secret proposal was incomplete; I requested "3 rope" but omitted the other items (screw and frame) that I agreed to leave for the other player. Although the rules state proposals must be complementary, submitting a proposal that only lists the items I want without explicitly accounting for the full set likely caused the system to reject the agreement or interpret it as a conflict. Next time, after reaching a verbal agreement, I will ensure my secret proposal explicitly lists the exact items I am claiming (e.g., "{[}Proposal: 3 rope{]}") while ensuring the other player's proposal covers the remainder. However, given the specific failure, I suspect the system requires the proposal to be a complete partition or I misinterpreted the "complementary" check. To be safe, I will strictly adhere to the agreed split: if I want the ropes, I will submit "{[}Proposal: 3 rope{]}" and trust the other player to submit the rest, but if the game requires me to define the full distribution in my proposal to ensure complementarity, I will format it as "{[}Proposal: 3 rope, 0 screw, 0 frame{]}" or simply ensure the other player's proposal is confirmed before I submit. Actually, looking closer at the failure, the most likely error is that the system expects the proposal to be a valid partition of the *entire* set or that my proposal was syntactically insufficient to prove complementarity against the other player's potential proposal. I will refine my strategy to: 1) Confirm the exact split verbally. 2) Submit a proposal that clearly lists only the items I am taking, but ensure the count matches the total available (e.g., "3 rope" is correct for ropes, but I must ensure the other player's proposal is "1 screw, 1 frame"). If the game logic requires me to explicitly state I am *not* taking the others to prove complementarity, I will adjust. But the most direct fix is to ensure the proposal format is exactly as agreed: "{[}Proposal: 3 rope{]}" is likely correct, so the error might be that I didn't wait for the other player to confirm their proposal or the system requires a specific format for the *entire* set. Wait, the rule says "If your proposals are complementary...". If I submit "3 rope" and the other submits "1 screw, 1 frame", they are complementary. The failure might be that I didn't actually submit the proposal correctly or the game ended prematurely. Let's assume the issue was that I didn't explicitly state the *full* allocation in my proposal to avoid ambiguity. I will try submitting a proposal that lists the items I want, but if that fails, I will try listing the items I *don't* want to ensure the system sees the complement. No, the rule says "indicating the items you want". The most likely mistake is that I didn't verify the other player's} \\
\\ 
\\ 
\texttt{You are playing a negotiation game in which you have to agree on how to divide a set of items among you and another player.} \\
\\ 
\texttt{Rules:} \\
\texttt{(a) You and the other player are given a set of items. Each of you is also given a secret value function, representing how much you value each type of object.} \\
\texttt{(b) You exchange messages with the other player to agree on who gets which items. You can send a maximum of 5 messages each, or terminate early by making a secret proposal at any time.} \\
\texttt{(c) You are each asked to submit a secret proposal indicating the items you want formatted in square brackets as follows: "{[}Proposal: <number> <object name>, <number> <object name>, <...>{]}"} \\
\texttt{(d) If your proposals are complementary, i.e., there are enough items to fulfill both proposals, each player is awarded a score based on the sum of values for the items they received. Otherwise, both of you get zero points.} \\
\texttt{(e) Your goal is to maximize the score you receive.} \\
\\ 
\texttt{Let us start.} \\
\\ 
\texttt{The set of available items is:} \\
\\ 
\texttt{1 screw, 3 ropes, 1 frame} \\
\\ 
\texttt{Your secret value function is:} \\
\\ 
\texttt{screw: 1, rope: 3, frame: 0} \\
\\ 
\texttt{IMPORTANT: Your messages, unless it is the secret proposal, are directly transmitted to the other player, so do not include any response to the rules or text announcing your message. To make a secret proposal, use the indicated format. Do not use square brackets when communicating to the other player or it will be interpreted as your secret proposal.} \\
            }
        }
    }
    & \\ \\

    \theutterance \stepcounter{utterance}  
    & \multicolumn{4}{p{0.6\linewidth}}{
        \cellcolor[rgb]{0.9,0.9,0.9}{
            \makecell[{{p{\linewidth}}}]{
                \texttt{\tiny{[P1$\rangle$GM]}}
                \texttt{Hello! I see we have 1 screw, 3 ropes, and 1 frame to divide.} \\
\\ 
\texttt{Based on my values, I place a high priority on the ropes (value 3 each) and have no use for the frame (value 0). The screw is worth 1 to me.} \\
\\ 
\texttt{To maximize our total score, it seems efficient if I take the 3 ropes. That would leave the 1 screw and 1 frame for you. Does that split work for you? If so, I can submit my proposal immediately.} \\
            }
        }
    }
    & & \\ \\

    \theutterance \stepcounter{utterance}  
    & & & \multicolumn{2}{p{0.3\linewidth}}{
        \cellcolor[rgb]{0.9,0.9,0.9}{
            \makecell[{{p{\linewidth}}}]{
                \texttt{\tiny{[GM$|$GM]}}
                \texttt{continue} \\
            }
        }
    }
    & & \\ \\

    \theutterance \stepcounter{utterance}  
    & & & \multicolumn{2}{p{0.3\linewidth}}{
        \cellcolor[rgb]{0.9,0.9,0.9}{
            \makecell[{{p{\linewidth}}}]{
                \texttt{\tiny{[GM$|$GM]}}
                \texttt{Hello! I see we have 1 screw, 3 ropes, and 1 frame to divide.} \\
\\ 
\texttt{Based on my values, I place a high priority on the ropes (value 3 each) and have no use for the frame (value 0). The screw is worth 1 to me.} \\
\\ 
\texttt{To maximize our total score, it seems efficient if I take the 3 ropes. That would leave the 1 screw and 1 frame for you. Does that split work for you? If so, I can submit my proposal immediately.} \\
            }
        }
    }
    & & \\ \\

    \theutterance \stepcounter{utterance}  
    & & \multicolumn{4}{p{0.6\linewidth}}{
        \cellcolor[rgb]{0.9,0.9,0.9}{
            \makecell[{{p{\linewidth}}}]{
                \texttt{\tiny{[GM$\rangle$P2]}}
                \texttt{You have attempted this same task before and did not succeed. These are the plans you wrote after those attempts.} \\
\texttt{Trial \#1: New plan: I failed because I submitted a proposal for only the 3 ropes, ignoring the frame and screw I verbally agreed to take. This left the other player with no items in their proposal (or a conflicting one), causing the proposals to be non{-}complementary and resulting in zero points. Next time, after negotiating an agreement, I will ensure my secret proposal explicitly lists *all* items I agreed to receive, formatted correctly. Specifically, I will submit "{[}Proposal: 3 rope, 1 frame, 1 screw{]}" to match the full set of items and ensure the proposals are complementary.} \\
\\ 
\\ 
\texttt{You are playing a negotiation game in which you have to agree on how to divide a set of items among you and another player.} \\
\\ 
\texttt{Rules:} \\
\texttt{(a) You and the other player are given a set of items. Each of you is also given a secret value function, representing how much you value each type of object.} \\
\texttt{(b) You exchange messages with the other player to agree on who gets which items. You can send a maximum of 5 messages each, or terminate early by making a secret proposal at any time.} \\
\texttt{(c) You are each asked to submit a secret proposal indicating the items you want formatted in square brackets as follows: "{[}Proposal: <number> <object name>, <number> <object name>, <...>{]}"} \\
\texttt{(d) If your proposals are complementary, i.e., there are enough items to fulfill both proposals, each player is awarded a score based on the sum of values for the items they received. Otherwise, both of you get zero points.} \\
\texttt{(e) Your goal is to maximize the score you receive.} \\
\\ 
\texttt{Let us start.} \\
\\ 
\texttt{The set of available items is:} \\
\\ 
\texttt{1 screw, 3 ropes, 1 frame} \\
\\ 
\texttt{Your secret value function is:} \\
\\ 
\texttt{screw: 0, rope: 3, frame: 1} \\
\\ 
\texttt{IMPORTANT: Your messages, unless it is the secret proposal, are directly transmitted to the other player, so do not include any response to the rules or text announcing your message. To make a secret proposal, use the indicated format. Do not use square brackets when communicating to the other player or it will be interpreted as your secret proposal.} \\
\\ 
\\ 
\texttt{Hello! I see we have 1 screw, 3 ropes, and 1 frame to divide.} \\
\\ 
\texttt{Based on my values, I place a high priority on the ropes (value 3 each) and have no use for the frame (value 0). The screw is worth 1 to me.} \\
\\ 
\texttt{To maximize our total score, it seems efficient if I take the 3 ropes. That would leave the 1 screw and 1 frame for you. Does that split work for you? If so, I can submit my proposal immediately.} \\
            }
        }
    }
    & \\ \\

    \theutterance \stepcounter{utterance}  
    & & & \multicolumn{4}{p{0.6\linewidth}}{
        \cellcolor[rgb]{0.9,0.9,0.9}{
            \makecell[{{p{\linewidth}}}]{
                \texttt{\tiny{[GM$\langle$P2]}}
                \texttt{That split works perfectly for me! I will take the 1 screw and 1 frame, and you can have the 3 ropes.} \\
\\ 
\texttt{{[}Proposal: 3 rope{]}} \\
            }
        }
    }
     \\ \\

    \theutterance \stepcounter{utterance}  
    & & & \multicolumn{2}{p{0.3\linewidth}}{
        \cellcolor[rgb]{0.9,0.9,0.9}{
            \makecell[{{p{\linewidth}}}]{
                \texttt{\tiny{[GM$|$GM]}}
                \texttt{proposal} \\
            }
        }
    }
    & & \\ \\

    \theutterance \stepcounter{utterance}  
    & & & \multicolumn{2}{p{0.3\linewidth}}{
        \cellcolor[rgb]{0.9,0.9,0.9}{
            \makecell[{{p{\linewidth}}}]{
                \texttt{\tiny{[GM$|$GM]}}
                [0, 3, 0]
            }
        }
    }
    & & \\ \\

    \theutterance \stepcounter{utterance}  
    & & \multicolumn{4}{p{0.6\linewidth}}{
        \cellcolor[rgb]{0.9,0.9,0.9}{
            \makecell[{{p{\linewidth}}}]{
                \texttt{\tiny{[P1$\langle$GM]}}
                \texttt{The other player has submitted a secret proposal. It is now your turn to submit a proposal of your own.} \\
            }
        }
    }
    & \\ \\

    \theutterance \stepcounter{utterance}  
    & \multicolumn{4}{p{0.6\linewidth}}{
        \cellcolor[rgb]{0.9,0.9,0.9}{
            \makecell[{{p{\linewidth}}}]{
                \texttt{\tiny{[P1$\rangle$GM]}}
                \texttt{{[}Proposal: 3 rope{]}} \\
            }
        }
    }
    & & \\ \\

    \theutterance \stepcounter{utterance}  
    & & & \multicolumn{2}{p{0.3\linewidth}}{
        \cellcolor[rgb]{0.9,0.9,0.9}{
            \makecell[{{p{\linewidth}}}]{
                \texttt{\tiny{[GM$|$GM]}}
                \texttt{proposal} \\
            }
        }
    }
    & & \\ \\

    \theutterance \stepcounter{utterance}  
    & & & \multicolumn{2}{p{0.3\linewidth}}{
        \cellcolor[rgb]{0.9,0.9,0.9}{
            \makecell[{{p{\linewidth}}}]{
                \texttt{\tiny{[GM$|$GM]}}
                [0, 3, 0]
            }
        }
    }
    & & \\ \\

    \theutterance \stepcounter{utterance}  
    & & & \multicolumn{2}{p{0.3\linewidth}}{
        \cellcolor[rgb]{0.9,0.9,0.9}{
            \makecell[{{p{\linewidth}}}]{
                \texttt{\tiny{[GM$|$GM]}}
                [[0, 3, 0], [0, 3, 0]]
            }
        }
    }
    & & \\ \\

    \theutterance \stepcounter{utterance}  
    & \multicolumn{4}{p{0.6\linewidth}}{
        \cellcolor[rgb]{0.95,0.95,0.95}{
            \makecell[{{p{\linewidth}}}]{
                \texttt{\tiny{[P2$\rightarrow$P1]}}
                \texttt{Post{-}episode Reflexion (not sent to GM):} \\
\texttt{New plan: I failed because I submitted my secret proposal immediately after the other player did, without first negotiating a confirmed agreement on the split. The other player's secret proposal likely requested items I needed (e.g., some ropes) or a split that didn't match my assumption, causing the proposals to be non{-}complementary. My strategy of assuming a split based on my own values without verifying the other player's needs was flawed. Next time, I will: 1) Explicitly ask the other player what items they want before making any assumptions. 2) Propose a specific split (e.g., "I take 3 ropes, you take 1 screw and 1 frame") and wait for their explicit verbal confirmation. 3) Only submit my secret proposal *after* they have confirmed they will submit the complementary proposal. If they do not confirm, I will negotiate further or adjust my offer to ensure the proposals will match. I will not submit a secret proposal until the verbal agreement is solidified.} \\
            }
        }
    }
    & & \\ \\

    \theutterance \stepcounter{utterance}  
    & & & \multicolumn{4}{p{0.6\linewidth}}{
        \cellcolor[rgb]{0.95,0.95,0.95}{
            \makecell[{{p{\linewidth}}}]{
                \texttt{\tiny{[P1$\rightarrow$P2]}}
                \texttt{Post{-}episode Reflexion (not sent to GM):} \\
\texttt{New plan: I failed because I submitted a proposal for only the 3 ropes, ignoring the frame and screw I verbally agreed to take. This left the other player with no items in their proposal (or a conflicting one), causing the proposals to be non{-}complementary and resulting in zero points. Next time, after negotiating an agreement, I will ensure my secret proposal explicitly lists *all* items I agreed to receive, formatted correctly. Specifically, I will submit "{[}Proposal: 3 rope, 1 frame, 1 screw{]}" to match the full set of items and ensure the proposals are complementary.} \\
            }
        }
    }
     \\ \\

\end{supertabular}
}

\end{document}
